
\begin{comment}

% SMT solvers orchestrating theory specific solvers around a boolean SAT solver.
% It is essentially ubiquitous that one of these theory specific solvers is the
% theory of uninterpreted functions, backed in some form by an \egraph. One of the
% possible mechanisms by which to treat universal quantifiers is e-matching.
% Universally quantified formulas may be annotated by trigger multi-patterns,
% which search the \egraph and instantiate the axiom upon a match.

% In distinction to egglog, SMT solvers are heavily backtracking based.

% At a high level, SMT solvers contain the pieces from which one might construct a
% solver similar to Egglog, but they are not arranged in this way.

% SMT solvers are difficult to debug and understand. The input language is highly
% declarative, with very little operational intent. Egglog's operational semantics
% is quite simple and predictable enough to be programmable.

% There is more content in the result returned by Egglog compared to a SMT solver.
% A SMT solver may either return unsat with a proof or SAT with a model. The
% egglog database only contains facts that can be justified by the application of
% rules. A SMT model has no requirements for justification.

% You can take an inferred fact from the egglog database, assert its negation to
% an SMT solver, and the SMT solver may return unsat (or unknown or loop forever).
% In this case, the egglog justification and the SMT unsat proof may align.
% However, there appears to be no obvious modality of use of an SMT solver that
% would find this derived fact and return it to the user.

% Z3 exposes it's rewriting simplifier via a simplify command. This simplifier
% however is not programmable, and does not use assumptions in the solver context
% that may be necessary, such as no-aliasing conditions in array expressions.

% Horn clauses specified in SMTLib are capable of expressing a subset of the
% syntax of egglog rules.

% Z3 also has a built in Datalog engine. This Datalog engine is largely disjoint
% from the main capabilities of Z3 and is mostly an entirely separate project in
% the same codebase.


\yz{what's a canonical reference for prolog?}

% TODO: mention unification variables cannot be removed

% z3 is not minimal. A solution to equality saturation will union everything together.

\subsection{E-graphs}

\Egraphs are first introduced in Greg Nelson's seminal thesis~\citep{nelson}
 in late 1970s 
 as a way of effectively deciding the theory of equalities.
A more efficient algorithm is later introduced by \citet{tarjan-congruence} 
 and the time complexity of this algorithm is analyzed.
\Egraphs are then used at the core of 
 various theorem provers and solvers~\citep{simplify, z3, cvc4}.
In the 2000s, 
 \egraphs are repurposed for program optimization \citet{eqsat,denali}.
The technique, known as equality saturation, 
 repeatedly performs non-destructive rewriting on the \egraphs 
 to grow a compact space of equivalent programs.
An extraction procedure is then applied to extract the optimal program.
In essence, equality saturation mitigates the phase-ordering problem by keeping 
 all programs.
This inspires later work on using \egraphs for
 translation validation~\citep{eqsat-llvm}, 
 floating-point arithmetic~\citep{herbie},
 semantic code search~\citep{yogo-pldi20},
 and computer-aided design~\citep{wu_siga19}.
However, a generic framework for \egraphs is not yet available,
 and developing \egraphs--based applications
 requires implementing \egraphs from scratch and is therefore a tedious effort.
Recently, a generic framework for \egraphs, 
 called \texttt{egg}, is developed~\citep{egg}.
As a result, many projects sprang up building domain-specific projects using \egraphs,
 including rewrite rule synthesis~\citep{ruler}, machine learning compiler~\citep{tensat,glenside},
 relational query optimization~\citep{spores}, and
 digital signal processing vectorization~\citep{diospyros}.

The connection between \egraphs and relational databases
 is first studied by \citet{relational-ematching}.
They propose to view an \egraph as a relational database,
 which makes \ematching orders of magnitude faster
 and prove desired theoretical properties.
However, to use this technique,
 one has to keep both the \egraph and its relational representation 
 and convert back and forth, which limits its practical adoptions.
We build on their work, 
 which only takes a static relational snapshot of \egraph each time,
 and explore how \egraphs as a relational database will behave dynamically.
This saves us from the labor of keeping and syncing between the two \egraph representations
 and further exploits the benefits of the relational \ematching approach.

Many works on improving \egraphs focus on the efficiency and usability of \egraphs.
% Some of these works are subsumed by the relational \egraphs.
Relational \egraphs bring a new perspective to some of these works.
For example, \citet{efficient-e-matching} studied
 the incremental maintenance problem of the \ematching procedure
 and one of its standard extension called multi-patterns.
\citet{tensat} also proposed an algorithm 
 to extend \egraph frameworks with multi-patterns for equality saturations.
In relational \egraphs, multi-patterns are supported naturally~\citep{relational-ematching},
 and incremental maintenance can be achieved using semi-na\"ive evaluation,
 which is a standard technique for evaluating Datalog programs~\citep{seminaive}.
Many applications extend the expressiveness of \egraphs 
 by writing domain-specific analyses in a general purpose language.
For example, to reason about lambda calculus in \egg, 
 a user may want to implement an analysis that
 manually tracks the set of free and bound variables~\citep{egg}
 in Rust.
Such analyses can be written as rules in pure relational \egraphs.
For some other works, 
 the problem of finding and understanding its relational dual 
 is still open to future research.
For instance,
 applications like SMT solvers not only want to know if two terms are equivalent, 
 but also why they are equivalent.
Techniques are developed to generate proofs for equivalences in \egraphs~\citep{pp-congr}\%
\revision{, a future work is to explore proof generation for \egglog programs}.
\delete{
 It is speculated that proofs for congruence closure in relational form 
 may just be database provenance~\citep{prov-semiring,prov-souffle}.
}
Moreover,
 scheduling is a critical component of equality saturation~\citep{egg}, 
 and a good scheduling algorithm is a key enabler of scalable equality saturations.
However, 
 for relational languages like Datalog,
 the scheduling problem is less studied,
 because Datalog programs are usually run to fixpoint,
 while the fixpoint for equality saturation is usually infinitary.
A future direction is to study the scheduling problem for relational \egraphs.

Other data structures for compact program representation 
 have also been long studied in the literature for decades,
 in particular version space algebras (VSAs)~\citep{vsa,flashmeta} 
 and finite tree automata~\citep{blaze, dace}.
Recently, it is shown that both \egraphs and VSAs are special cases of finite tree automata~\citep{vsa-eq-fta}.
A natural question is therefore if we can encode 
 VSAs, finite tree automata, and operations over them in a relational approach,
 and how we can possibly benefit from such an encoding for tasks 
 like program synthesis and program optimization.

\subsection{Relational databases and Datalog}

Our work is closely related to works on Datalog and relational database.
Relational \egraphs are directly inspired by work on data dependencies and the chase.
Equality generating dependencies (EGDs) are data dependencies of the form 
 $\forall \vec{x}.\lambda(\vec{x}) \rightarrow x_i=x_j$,
 which asserts the equality between terms under conditions given by predicate $\lambda$.
EGDs generalizes the congruences in equality saturation, 
 e.g., the congruence property of binary operator \textit{add} can be written as
 $\forall x,y,z_1,z_2.\textit{add}(x, y, z_1), \textit{add}(x, y, z_2)\rightarrow z_1=z_2$.
Tuple generating dependencies (TGDs) are another kind of data dependencies of the form
 $\forall \vec{x}, \lambda(\vec{x})\rightarrow \exists \vec{y}, \rho(\vec{x}, \vec{y})$,
 which essentially generalizes Datalog rules with existential quantifiers in the body.
TGDs generalizes rewrite rules in equality saturation.
For example, the associativity is expressed as TGD
 $\forall x,y,xy,z,xyz. \textit{add}(x, y, xy), \textit{add}(xy, z, xyz)\rightarrow \exists yz.\textit{add}(y, z, yz), \textit{add}(x, yz, xyz)$.
The chase is a family of iterative algorithms
 for reasoning about data dependencies~\citep{chase-revisited, bench-chase}.
Therefore,
 equality saturation can be effectively viewed as a chase algorithm.
As a chase procedure, equality saturation has many interesting properties that are worth further study:
While there may be many finite universal models to data dependencies in the chase,
 in equality saturation,
 there will be at most one finite universal model, 
 which is the core.
Moreover,
 equality saturation terminates 
 if and only if there is a finite universal model of the given dependencies, 
 which equality saturation will output,
In contrast,
 the (non-core) chase does not necessarily terminate 
 even when a finite universal solution of the input dependencies exists,
 or such a finite solution is very expensive to compute (with the core chase~\citep{chase-revisited}).
A future direction is 
 to understand equality saturation using the theories developed 
 in database research.

Many applications using \egraphs rely on domain-specific analyses.
To express these analyses in a generic \egraph framework, 
 \egg proposed a framework called \eclass analyses.
%  which maintains a lattice associated with each \eclass monotonically.
An \eclass analysis can be thought of as an aggregation of 
 information in the corresponding \eclass's \enodes and their children's \eclass analyses.
\Eclass analyses are interdependent,
 because the \eclass analysis of one \eclass depends on its children's \eclass analyses.
Therefore,
 \eclass analyses can be formulated as recursive aggregates in Datalog.
However, unconstrained aggregates within recursions are dangerous, 
 as it can lead to non-monotonic programs and there may not be (stable) models,
 so in practice aggregate stratification is enforced.
Several approaches are proposed to loosen the aggregate stratification requirement 
~\citep{mono-agg,bloom-lattice,datalogo,prov-semiring,flix}.
Among these approaches,
 the \eclass analyses naturally matches the lattice semantics proposed by Flix~\citep{flix},
 as each \eclass analysis monotonically maintains a lattice, which is associated with each \eclass,
In Flix, relations are optionally annotated with a lattice,
 which aggregates over values passed it according to the lattice join operation.
% This is precisely the semantics of \eclass analyses in relational form.
Our relational \egraphs extends Flix by allowing multiple atoms in the head.
A similar lattice-based approach is studied 
 in the Bloom$^\text{L}$ language~\citep{bloom-lattice}.

Following relational \ematching~\citep{relational-ematching}, 
 we use worst-case optimal join for solving queries over relational \egraphs.
Different from relational \ematching,
 which only needs to consider static snapshots of a relational database
 (i.e., no writes),
 we have to also consider insertions and updates to the database.
Several previous works considered 
 adopting worst-case optimal joins for practical systems.
Two critical dimensions in the design space is the design of indexes
 and query planning.
For indices,
 EmptyHeaded only considers scenarios with static dataset~\citep{emptyheaded}
 and therefore their indexes are read-optimized and need to be precomputed.
Our context for using worst-case optimal join requires to update the database
 and is therefore more similar to LogicBlox.
LogicBlox is a commercial database system that uses leapfrog triejoin (LFTJ)~\citep{lftj}, 
 which is worst-case optimal,
 for its query processing~\citep{logicblox}.
To support both LFTJ and efficient updates, 
 LogicBlox uses write-optimized B-trees for indexes.
Different indexes are created and maintained 
 when they are used by query plans generated from the query optimizer.
Recently, 
 \citet{umbra-wcoj} implements a worst-case optimal join 
 inside the Umbra database system.
In their design,
 the indices are built on-the-fly during join processing
 and are optimized for fast building.
This design allows efficient updates to databases 
 since no additional indices need to be maintained during updates.

In terms of query optimization,
 LogicBlox uses a sampling-based technique to pick a good query plan.
EmptyHeaded uses generalized hypertree decomposition (GHD),
 which allows it to provide guarantees even stronger than 
 those provided by standard worst-case optimal joins.
\citep{umbra-wcoj} uses cardinality information to optimize its query plans,
 and allows hybrid query plans that use binary joins like hash joins
 when it is deemed that worst-case optimal join does not offer a benefit.
Currently, the query optimizer of our relational \egraph is relatively simple,
 and we hope to study and adopt techniques used in other database systems 
 that use worst-case optimal joins.

\subsection{Egg}
Egglog is a generalization and homogenization of the work in Egg. A novel
feature of Egg compared to previous work on egraphs was the presence of
Analyses, data annotated on top of eclasses. Egglog offers a concrete
s-expression based syntax that avoids the need to code and compile Rust, a new
high performance relational  e-matching engine, and allows easy and
compositional creation of analyses in the datalog style.

\subsection{Datalog}
Datalog is a variant of prolog that is commonly executed in a bottom up fashion
using database algorithms. In its narrow definition, it does not allow for
structured data or unbounded arithmetic operations. These restrictions allow for
trivial arguments for it's termination. The final database is bounded above by a
fully saturated database filled with every combination of initial symbols.

There are a number of enhancements that one can make to Datalog.

Datalog can include terms. If new terms are constructable, the simple
termination guarantees of Datalog are out the window.

Some extensions to Datalog can be interpreted as logically motivated ways to
achieve deletion.

Lattice types are one possible addition to Datalog. The canonical example of
such a Datalog is Flix. The axioms of a lattice give the leeway necessary to
unconstrain the order and multiplicity of application of the datalog rules. The
resolution of collisions between distinct lattice values is very reminiscent of
egglog's resolution of congruence violations.

Another feature which achieves a principled deletion is that of subsumption, a
feature that souffle has. One may declare conditions upon which one entry in the
database is subsumed by another.

\subsection{The Chase}
Egglog's operational semantics are some variant of the chase. Indeed Datalog
itself can be seen as a variant of the chase.

\pz{Are there restrictions of egglog compared to full chase that make it 
simpler/faster/tuned to egraph use cases? Or is it just emphasis? 
In the standard chase, the functional nature of a relation would be 
described via a clause, whereas it is built in facet of the definition of a 
function in egglog.}

Datalog+/- https://www.cs.ox.ac.uk/publications/publication6640-abstract.html
https://www.aaai.org/ocs/index.php/KR/KR14/paper/viewFile/7965/7972 is a datalog
language that allows equality and existentials in the heads of it's rules.
Egglog is perhaps an example of a Datalog+/-. It is not clear that Datalog+/- is
more than a theoretical exercise.

\subsection{Prolog}
Prolog is the canonical logic programming language.
https://arxiv.org/abs/2201.10816 Prolog is a general purpose programming
language and includes many non-monotonic operational constructs such as cut and
IO.

Prolog is distinguished by having first class support for unification variables
and backtracking search. As a pure logical inference system, it has it's
deficits. Prolog by default has an incomplete search strategy and can easily get
into infinite loops.

The elegant combination of Horn clauses with a union-find suggests that egglog
might be an answer to how to add unification variables to datalog. This is
however not the case. The "unification variables" in egglog are global, whereas
prolog's are branch local and backtrackable.

Prolog terms are compared to each other via structural unification. This is very
different from egglog, in which structural dissimilar terms can indeed be equal
to each other.

Constraint Logic Programming and Constrained Horn clauses are two related
languages in which horn clauses additionally allow theory specific constraints.
There are disparate methods for solving these systems, but a typical one is a
prolog like search with a constraint store. Egglog could be said to be
expressing constrained horn clauses, where the theory is the theory of
uninterpreted functions. This theory however is deeply coupled to the terms of
the language of egglog, in a manner that is not typically implied by the two
formalisms.

\subsection{Constraint Handling Rules}
Constraint Handling Rules are a multiset term rewriting system. Because they are
multisets, they have a database character. With deduplication rules, it is
simple to embed datalog into CHR. Constraint Handling Rules have been used to
implement union-finds, and term rewriting systems. They are a very flexible
system for defining propagating information about bespoke constraints. However,
it is known to be the case that solvers developed specifically for individual
systems can outperform the general purpose mechanisms of CHR. It is possible to
embed e-matching and congruence closure into CHR, but again at a performance
cost.

\subsection{Proto-Egglogs}
Egg-lite is an e-graph implementation based around SQLite.

egglog0 is a datalog built upon the egg library using a simple e-matching
multipattern implementation. It had no support for primitives or for analyses,
and used a weaker non relational e-matching engine.

souffle-egglog is an embedding of e-graphs into Souffle datalog. Ultimately, the
lack of rebuild as an intrinsic operation makes this encoding too inefficient.
It is also highly burdensome to perform manual flattening of terms in rules
without at least a metaprogramming system.

\subsection{Formulog}
Formulog is a system that combines datalog with SMT solvers by adding SMT
expressions as a primitive type and SMT queries as a primitive function to a
datalog engine. Egglog and Formulog have some overlapping capabilities, however
egglog deeply integrates the SMT-internal egraph into it's structure, whereas
Formulog allows for search like backtracking disjunctional behavior when it
hands over control to the SMT solver.

An somewhat unfair set of equations
$$
formulog = Datalog + smt
$$
$$
smt = egraph + sat
$$
$$
egglog = egraph + datalog = formulog - sat
$$
\subsection{Satisfiability Modulo Theories (SMT)}
SMT solvers orchestrating theory specific solvers around a boolean SAT solver.
It is essentially ubiquitous that one of these theory specific solvers is the
theory of uninterpreted functions, backed in some form by an \egraph. One of the
possible mechanisms by which to treat universal quantifiers is e-matching.
Universally quantified formulas may be annotated by trigger multi-patterns,
which search the \egraph and instantiate the axiom upon a match.

In distinction to egglog, SMT solvers are heavily backtracking based.

At a high level, SMT solvers contain the pieces from which one might construct a
solver similar to Egglog, but they are not arranged in this way.

SMT solvers are difficult to debug and understand. The input language is highly
declarative, with very little operational intent. Egglog's operational semantics
is quite simple and predictable enough to be programmable.

There is more content in the result returned by Egglog compared to a SMT solver.
A SMT solver may either return unsat with a proof or SAT with a model. The
egglog database only contains facts that can be justified by the application of
rules. A SMT model has no requirements for justification.

You can take an inferred fact from the egglog database, assert its negation to
an SMT solver, and the SMT solver may return unsat (or unknown or loop forever).
In this case, the egglog justification and the SMT unsat proof may align.
However, there appears to be no obvious modality of use of an SMT solver that
would find this derived fact and return it to the user.

Z3 exposes it's rewriting simplifier via a simplify command. This simplifier
however is not programmable, and does not use assumptions in the solver context
that may be necessary, such as no-aliasing conditions in array expressions.

Horn clauses specified in SMTLib are capable of expressing a subset of the
syntax of egglog rules.

Z3 also has a built in Datalog engine. This datalog engine is largely disjoint
from the main capabilities of Z3 and is mostly an entirely separate project in
the same codebase.

\subsection{Algebraic Theories}
Algebraic theories are those specified by universally quantified equational
axioms, such as groups or rings. Egraph rewriting is a natural theorem proving
technique to apply to these axiom systems. Three related systems that extend
algebraic theories are essentially algebraic theories, generalized algebraic
theories
\href{https://www.sciencedirect.com/science/article/pii/0168007286900539?via%3Dihub} , and partial Horn Logic \href{http://www2.math.uu.se/~palmgren/partialalgebras_pre.pdf} . These system allow partial operators and axioms whose well formedness is conditioned upon the truth of other equations, or well typedness conditions. The pure sort fragment of egglog considered as a logic is very similar to these systems and can be used to encode a prover for them. These systems may form a theoretical justification for the design of egglog.

\end{comment}